% GENERATED FILE. Edit canonical lessons, then run npm run generate:books.
% canonical-source-hash: fnv1a64:5cbf6efd1410e989
% canonical-lessons: ML-C186-pankuvekkuka, ML-C186-kettippitikkuka, ML-C186-valakkituka, ML-C186-pinannuka, ML-C186-kaliyakkuka

\chapter{To share, To hug, To quarrel, To sulk, To tease}
\label{ch:ml-to-share-to-hug-to-quarrel-to-sulk-to-tease}
\hlchaptermodality{186}

\begin{chapteropening}
\textbf{By the end of this chapter:} \emph{I can say to share, to hug, to quarrel, to sulk and to tease.}

Complete the last lesson of chapter 186: I can say to share, to hug, to quarrel, to sulk and to tease.
\end{chapteropening}

\section[\texorpdfstring{\ml{പങ്കുവെക്കുക} (paṅkuvekkuka)}{പങ്കുവെക്കുക (paṅkuvekkuka)}]{\ml{പങ്കുവെക്കുക} (paṅkuvekkuka) — to share}

\label{lesson:ML-C186-pankuvekkuka}



\begin{quote}
Before the new one: say the Malayalam for to fly, then the Malayalam for to follow.
\end{quote}

\subsection*{You'll want to know: \ml{പങ്കുവെക്കുക}}
\textbf{\ml{പങ്കുവെക്കുക}} — \emph{paṅkuvekkuka} — \textquotedblleft{}to share\textquotedblright{}.

\begin{grammarlens}[title={What you've built}]
One more word for this chapter.
\end{grammarlens}

\subsection*{Your turn}
\begin{itemize}
\raggedright
  \item \emph{Say it:} \emph{paṅkuvekkuka}
  \item \emph{Say it:} \emph{paṅkuvekkuka}, once more
  \item \emph{Say it:} say \emph{paṅkuvekkuka}
  \item \emph{Recall:} say the Malayalam for to fly, then the Malayalam for to follow, then say \emph{paṅkuvekkuka} again
\end{itemize}

\begin{tcolorbox}[breakable,colback=teal!4,colframe=teal!35!black,title={Before you move on}]
What does \ml{പങ്കുവെക്കുക} mean? (\textquotedblleft{}To share\textquotedblright{}.) Say it once more.
\end{tcolorbox}
\section[\texorpdfstring{\ml{കെട്ടിപ്പിടിക്കുക} (keṭṭippiṭikkuka)}{കെട്ടിപ്പിടിക്കുക (keṭṭippiṭikkuka)}]{\ml{കെട്ടിപ്പിടിക്കുക} (keṭṭippiṭikkuka) — to hug}

\label{lesson:ML-C186-kettippitikkuka}



\begin{quote}
Before the new one: say the Malayalam for to follow, then the Malayalam for to share.
\end{quote}

\subsection*{You'll want to know: \ml{കെട്ടിപ്പിടിക്കുക}}
\textbf{\ml{കെട്ടിപ്പിടിക്കുക}} — \emph{keṭṭippiṭikkuka} — \textquotedblleft{}to hug\textquotedblright{}.

\begin{grammarlens}[title={What you've built}]
One more word for this chapter.
\end{grammarlens}

\subsection*{Your turn}
\begin{itemize}
\raggedright
  \item \emph{Say it:} \emph{keṭṭippiṭikkuka}
  \item \emph{Say it:} \emph{keṭṭippiṭikkuka}, once more
  \item \emph{Say it:} say \emph{keṭṭippiṭikkuka}
  \item \emph{Recall:} say the Malayalam for to follow, then the Malayalam for to share, then say \emph{keṭṭippiṭikkuka} again
\end{itemize}

\begin{tcolorbox}[breakable,colback=teal!4,colframe=teal!35!black,title={Before you move on}]
What does \ml{കെട്ടിപ്പിടിക്കുക} mean? (\textquotedblleft{}To hug\textquotedblright{}.) Say it once more.
\end{tcolorbox}
\section[\texorpdfstring{\ml{വഴക്കിടുക} (vaḻakkiṭuka)}{വഴക്കിടുക (vaḻakkiṭuka)}]{\ml{വഴക്കിടുക} (vaḻakkiṭuka) — to quarrel}

\label{lesson:ML-C186-valakkituka}



\begin{quote}
Before the new one: say the Malayalam for to share, then the Malayalam for to hug.
\end{quote}

\subsection*{You'll want to know: \ml{വഴക്കിടുക}}
\textbf{\ml{വഴക്കിടുക}} — \emph{vaḻakkiṭuka} — \textquotedblleft{}to quarrel\textquotedblright{}.

\begin{grammarlens}[title={What you've built}]
One more word for this chapter.
\end{grammarlens}

\subsection*{Your turn}
\begin{itemize}
\raggedright
  \item \emph{Say it:} \emph{vaḻakkiṭuka}
  \item \emph{Say it:} \emph{vaḻakkiṭuka}, once more
  \item \emph{Say it:} say \emph{vaḻakkiṭuka}
  \item \emph{Recall:} say the Malayalam for to share, then the Malayalam for to hug, then say \emph{vaḻakkiṭuka} again
\end{itemize}

\begin{tcolorbox}[breakable,colback=teal!4,colframe=teal!35!black,title={Before you move on}]
What does \ml{വഴക്കിടുക} mean? (\textquotedblleft{}To quarrel\textquotedblright{}.) Say it once more.
\end{tcolorbox}
\section[\texorpdfstring{\ml{പിണങ്ങുക} (piṇaṅṅuka)}{പിണങ്ങുക (piṇaṅṅuka)}]{\ml{പിണങ്ങുക} (piṇaṅṅuka) — to sulk, to fall out}

\label{lesson:ML-C186-pinannuka}



\begin{quote}
Before the new one: say the Malayalam for to hug, then the Malayalam for to quarrel.
\end{quote}

\subsection*{You'll want to know: \ml{പിണങ്ങുക}}
\textbf{\ml{പിണങ്ങുക}} — \emph{piṇaṅṅuka} — \textquotedblleft{}to sulk, to fall out\textquotedblright{}.

\begin{grammarlens}[title={What you've built}]
One more word for this chapter.
\end{grammarlens}

\subsection*{Your turn}
\begin{itemize}
\raggedright
  \item \emph{Say it:} \emph{piṇaṅṅuka}
  \item \emph{Say it:} \emph{piṇaṅṅuka}, once more
  \item \emph{Say it:} say \emph{piṇaṅṅuka}
  \item \emph{Recall:} say the Malayalam for to hug, then the Malayalam for to quarrel, then say \emph{piṇaṅṅuka} again
\end{itemize}

\begin{tcolorbox}[breakable,colback=teal!4,colframe=teal!35!black,title={Before you move on}]
What does \ml{പിണങ്ങുക} mean? (\textquotedblleft{}To sulk, to fall out\textquotedblright{}.) Say it once more.
\end{tcolorbox}
\section[\texorpdfstring{\ml{കളിയാക്കുക} (kaḷiyākkuka)}{കളിയാക്കുക (kaḷiyākkuka)}]{\ml{കളിയാക്കുക} (kaḷiyākkuka) — to tease, to make fun of}

\label{lesson:ML-C186-kaliyakkuka}



\begin{quote}
Before the new one: say the Malayalam for to quarrel, then the Malayalam for to sulk.
\end{quote}

\subsection*{You'll want to know: \ml{കളിയാക്കുക}}
\textbf{\ml{കളിയാക്കുക}} — \emph{kaḷiyākkuka} — \textquotedblleft{}to tease, to make fun of\textquotedblright{}.

\begin{grammarlens}[title={What you've built}]
One more word for this chapter.
\end{grammarlens}

\subsection*{Your turn}
\begin{itemize}
\raggedright
  \item \emph{Say it:} \emph{kaḷiyākkuka}
  \item \emph{Say it:} \emph{kaḷiyākkuka}, once more
  \item \emph{Say it:} say \emph{kaḷiyākkuka}
  \item \emph{Recall:} say the Malayalam for to quarrel, then the Malayalam for to sulk, then say \emph{kaḷiyākkuka} again
\end{itemize}

\begin{tcolorbox}[breakable,colback=teal!4,colframe=teal!35!black,title={Before you move on}]
What does \ml{കളിയാക്കുക} mean? (\textquotedblleft{}To tease, to make fun of\textquotedblright{}.) Say it once more.
\end{tcolorbox}
